\documentclass[11pt]{amsart}
\usepackage[T1]{fontenc}
\usepackage{lmodern}
\usepackage{amsmath,amssymb,amsthm}
\usepackage{booktabs}
\usepackage[margin=1.1in]{geometry}
\usepackage[hidelinks]{hyperref}
\hypersetup{
  pdftitle={The dimer entropy of the king's graph},
  pdfauthor={Dawid Ry\l{}ko},
  pdfsubject={MSC2020: Primary 05C70; Secondary 82B20, 82B80, 05A16},
  pdfkeywords={king's graph, domicule tilings, perfect matchings, dimer model, entropy,
    transfer matrix},
  pdfcreator={LaTeX with hyperref},
}
\usepackage{url}
\usepackage{microtype}
\emergencystretch=2em

\newtheorem{theorem}{Theorem}[section]
\newtheorem{lemma}[theorem]{Lemma}
\newtheorem{proposition}[theorem]{Proposition}
\newtheorem{conjecture}[theorem]{Conjecture}
\theoremstyle{definition}
\newtheorem{definition}[theorem]{Definition}
\newtheorem{problem}[theorem]{Open problem}
\theoremstyle{remark}
\newtheorem{remark}[theorem]{Remark}

\newcommand{\Z}{\mathbb{Z}}
\newcommand{\N}{\mathbb{N}}
\newcommand{\fl}{\operatorname{fl}}
\newcommand{\supp}{\operatorname{supp}}
\newcommand{\file}[1]{\texttt{#1}}
% Visible placeholder for anything that must be filled in before submission.
\newcommand{\TODO}[1]{\textbf{[TODO: #1]}}

\title{The dimer entropy of the king's graph}
\author{Dawid Ry\l{}ko}
\address{Independent researcher, Bielsko-Bia\l{}a, Poland}
\email{dawid.rylko@gmail.com}
\date{October 2026}
\thanks{This work is licensed under CC BY 4.0, \url{https://creativecommons.org/licenses/by/4.0/}.}
\subjclass[2020]{Primary 05C70; Secondary 82B20, 82B80, 05A16}
\keywords{king's graph, domicule tilings, perfect matchings, dimer model, entropy,
  transfer matrix}

\begin{document}

\begin{abstract}
We compute the entropy per site $s$ of the perfect matchings of the king's graph, the square
lattice with both diagonals; equivalently, of the domicule tilings counted in OEIS A239264. The
graph is not planar, and its boards with both sides at least $3$ admit no Pfaffian orientation, so
Kasteleyn's method does not apply. Row transfer matrices on cylinders of circumference up to $25$,
split by a conserved parity of the dimers crossing a loop, converge exponentially; a L\"uscher-type
extrapolation with predictive backtests and an error budget gives
\[
  s = 0.5655119110413253 \pm 6\cdot 10^{-16},
\]
a non-rigorous estimate. Computer-assisted Collatz--Wielandt certificates, combined with a bound
for the matchings across a cut that depends on the runs of crossing cells, prove
$0.545343749727 \le s \le 0.582767477484$. We also estimate the surface free energy of a free edge,
$\sigma = 0.242017935766251 \pm 9\cdot 10^{-15}$, the corner free energy, which gives the asymptotics of
A239273, and the decay length $\xi = 1.235 \pm 0.004$ of the splitting between the two parity
sectors. Two sign properties of the strip growth rates, convexity in the even width and a
one-sided bound on the odd increments, hold numerically and would give
$0.565511911035 \le s \le 0.565511911049$.
\end{abstract}

\maketitle

\section{Introduction}

A \emph{domicule} is a pair of cells of the square grid that share a side or a corner. A domicule
tiling of a $K \times L$ board is a partition of its cells into domicules. Two diagonal domicules
may cross inside a unit square; for example, the $2\times 2$ board has $3$ tilings, one of which
uses both diagonals. Equivalently, a domicule tiling is a
perfect matching of the \emph{king's graph} on the board, the graph whose edges join cells at
chessboard king distance $1$. The numbers $Z(K, L)$ of such tilings form the array
A239264 of the OEIS \cite{oeisA239264}, and $Z(2n, 2n)$ is A239273 \cite{oeisA239273}; related
entries such as A239265, A239268 and A243510 \cite{oeisA239265,oeisA239268,oeisA243510} record
further exact counts. None of these entries gives the growth constant, i.e.\ the entropy per
site
\[
  s = \lim_{K, L \to \infty} \frac{\log Z(K, L)}{KL}.
\]

For dominoes (the square lattice) $s$ is known exactly, $G/\pi$ with Catalan's constant $G$, by
the Pfaffian method of Kasteleyn and of Temperley and Fisher \cite{kasteleyn1961,
temperleyfisher1961}. Kasteleyn's method needs a Pfaffian orientation, and planar graphs have
one. The triangular lattice (one diagonal per square) is planar: Montroll gave an orientation
for it \cite{montroll1964}, its entropy is tabulated by Wu \cite{wu2006}, and Fendley, Moessner
and Sondhi computed its correlations \cite{fms2002}. The king's graph is not planar (the
$5 \times 5$ board already has $72 > 3\cdot 25 - 6$ edges), but planarity is only sufficient for a
Pfaffian orientation. We show that no king board with both sides at least $3$ has one
(Theorem~\ref{thm:pfaff}), so a single Pfaffian cannot count its domicule tilings, and no
determinantal formula for $Z$ is known. Wildeboer and Seidel \cite{wildeboerseidel2020} give a
``pre-Kasteleyn'' orientation of this lattice whose Pfaffian is the signed sum
$\sum_D (-1)^{n_c(D)}$ over the dimer coverings $D$, where $n_c(D)$ is the number of pairs of
crossing diagonal dimers, and they note that the unsigned problem is not addressable by Pfaffian
methods. Yao and Kivelson \cite{yaokivelson2012} write the variant without crossing dimers as a
Grassmann integral with a quartic interaction. Wildeboer, Nussinov and Seidel
\cite{wildeboernussinovseidel2020} and Wildeboer et al.\ \cite{wildeboer2026} solve different,
constrained dimer models on lattices with crossed plaquettes.

Sandvik and Moessner \cite{sandvikmoessner2006} (see also the earlier preprint \cite{sandvik2003})
simulated dimers on the nearest- and next-nearest-neighbour bonds of the square lattice by Monte
Carlo, at a small fugacity $w_2$ of the diagonal dimers between $e^{-6}$ and $e^{-3}$; they show
that the diagonal dimers deconfine monomers, with a correlation length $\propto 1/w_2$. The
equal-weight point $w_2 = 1$ studied here was not simulated, and no entropy is given. Finch's
errata \cite[\S 5.12]{finch2020errata} list a constant $1.34254\ldots 1.34265$ for the square
lattice with ``crossed diagonal bonds''; it is the growth constant $1.342643951124\ldots$ of
non-attacking kings (independent sets of the king's graph, OEIS A247413 \cite{oeisA247413}), not a
dimer constant. Series expansions for dimer and monomer--dimer problems on the standard two- and
three-dimensional lattices go back to Nagle \cite{nagle1966} and Gaunt \cite{gaunt1969}; Nagle's
expansion has the Bethe approximation as its zeroth order (Table~\ref{tab:context}). To our
knowledge no numerical value of $s$ has been published.

\emph{Related methods.} Calkin and Wilf \cite{calkinwilf1998} bounded the hard-square entropy by
eigenvalues of strip transfer matrices; Friedland and Peled \cite{friedlandpeled2005} developed
the approach for monomer--dimer systems, and Louidor and Marcus \cite{louidormarcus2010}
sharpened the lower bounds with Rayleigh quotients. These bounds converge fast, but their proofs
need a symmetric, reflection-positive transfer, which the king lattice does not have
(\S\ref{sec:conj}). Pavlov \cite{pavlov2012}, Marcus and Pavlov \cite{marcuspavlov2013} and
Gamarnik and Katz \cite{gamarnikkatz2009} compute entropies rigorously from strong spatial
mixing. For perfect matchings, the parity of the number of dimers crossing a closed curve is
fixed by the number of sites it encloses, a non-local constraint that these arguments do not
handle as they stand. Without rigour, corner transfer matrices gave the hard-square entropy to
$43$ decimals \cite{baxter1999}, and their renormalization-group variant
\cite{nishinookunishi1996} is a natural route to more digits here; the crossing diagonals and the
parity sector of \S\ref{sec:model} need a dedicated encoding, which we leave open. We use row
transfer matrices because they are exact, easy to check against brute force, and give the
estimate and the certificates from the same code.

\medskip\noindent\textbf{Results.}
\begin{enumerate}
\item No king board with both sides at least $3$ has a Pfaffian orientation
  (Theorem~\ref{thm:pfaff}).
\item $s = 0.5655119110413253 \pm 6\cdot10^{-16}$ (\S\ref{sec:value}), i.e.\ $2s = 1.13102382208265$
  per dimer and $e^{s} = 1.76034869404198$, both rounded to $14$ decimals. The error bar is not
  rigorous.
\item Rigorously, $0.545343749727 \le s \le 0.582767477484$ (Theorem~\ref{thm:bounds}); the upper bound
  uses the run structure of the matchings across a cut (Lemma~\ref{lem:runs}).
\item The surface free energy of a free edge is $\sigma = 0.242017935766251 \pm 9\cdot10^{-15}$, and
  the corner free energy is $\kappa = 0.36319625605 \pm 3\cdot10^{-10}$, which gives the
  asymptotics of A239273 (\S\ref{sec:surface}).
\item On even cylinders the two parity sectors split by $\Delta_w \approx e^{-w/\xi}$ with
  $\xi = 1.235 \pm 0.004$, and all finite-size corrections we see are exponentially small
  (\S\ref{sec:z2}), in contrast with the critical square lattice.
\item Convexity of $q \mapsto \log\lambda_{2q}$ (L$_\infty$) and a bound on the odd increments (O)
  of the strip growth rates hold wherever we can compute and would give
  $0.565511911035 \le s \le 0.565511911049$ (\S\ref{sec:conj}). The stronger log-convexity of the
  finite counts is not structural; a companion note \cite{rylko2026lnote} treats it.
\end{enumerate}
The same code reproduces the exactly known entropies of the square and triangular lattices
(\S\ref{sec:checks}). All scripts, outputs and the map from every number in this paper to the
file it comes from are in the repository (\S\ref{sec:repro}).

\section{Model, transfer matrix and parity sectors}\label{sec:model}

\subsection{Definitions}
The king's graph has vertex set $\Z^2$, and $(r, c) \sim (r', c')$ if
$\max(|r - r'|, |c - c'|) = 1$. The board $B(K, L) = \{0, \dots, K-1\} \times \{0, \dots, L-1\}$
has $K$ rows and $L$ columns; $Z(K, L)$ is the number of perfect matchings of the induced
subgraph ($Z = 0$ if $KL$ is odd), and $Z(K, L) = Z(L, K)$. The cylinder of circumference
$w \ge 3$ identifies column $c$ with $c + w$; its count is $Z_{\mathrm{cyl}}(K, w)$. Let
\[
  \bar s = \limsup \frac{\log Z(K, L)}{KL}, \qquad \underline s = \liminf \frac{\log Z(K, L)}{KL},
\]
over $K, L \to \infty$ with $KL$ even. The equality of the two limits is standard for dimer
models on $\Z^2$-periodic graphs (a sub- and superadditivity argument, as for the dimer problem on
$\Z^d$ \cite{friedlandpeled2005}), and for the king's graph it is shown below;
Theorem~\ref{thm:bounds} bounds $\underline s$ and $\bar s$ separately and does not use it.

For a strip of width $k$, stacking boards gives $Z(2i + 2j, k) \ge Z(2i, k) Z(2j, k)$, so by
Fekete's lemma \cite{fekete1923} the limit
\begin{equation}\label{eq:lambda}
  \lambda_k^2 = \lim_{j \to \infty} Z(2j, k)^{1/j} = \sup_j Z(2j, k)^{1/j}
\end{equation}
exists. We call $\lambda_k$ the growth rate of the free strip and write $a_k = \log \lambda_k$.
Placing strips side by side gives $\lambda_{k_1 + k_2} \ge \lambda_{k_1}\lambda_{k_2}$, so
$a_k/k \to s^* = \sup_k a_k/k$. By \eqref{eq:lambda}, $Z(2j, L) \le \lambda_L^{2j}$, so
$\log Z(2j, L)/(2jL) \le a_L/L \le s^*$; for odd $K$ (then $L$ is even) $Z(K, L) \le Z(K+1, L)$
(add a row of horizontal dominoes). Hence $\bar s \le s^*$, and with $\underline s \ge s^*$ (the
proof of the lower bound in \S\ref{sec:bounds}) the limit $s$ exists and equals
$s^* = \lim_k a_k/k$.

\subsection{Transfer matrix}
We scan a strip of width $w$ cell by cell in row-major order. A cell is either already covered,
or it is matched to one of its later neighbours: right $(r, c+1)$, down-left $(r+1, c-1)$, down
$(r+1, c)$ or down-right $(r+1, c+1)$, if that neighbour exists and is free. The state records
the occupancy of the next $w + 2$ cells; on the cylinder the neighbours are taken modulo $w$ and
one extra bit records the cell $(r+1, w-1)$ claimed by $(r, 0)$ through the wrap-around
down-left edge. A row is $w$ cell steps, $T = T_{w-1}\cdots T_0$, and
\[
  Z(K, w) = e_0^{\top} T^K e_0,
\]
with $e_0$ the empty state. Between rows only the $w$ bits of the next row can be set, so $T$
acts on $2^w$ states: the set $A \subseteq \{0, \dots, w-1\}$ of cells of the next row already
covered from the current one.

\subsection{Parity sectors}
Let $A_r$ be the set of cells of row $r$ covered from row $r - 1$ and $h_r$ the number of
dimers inside row $r$. Counting the cells of row $r$ gives $w = |A_r| + 2h_r + |A_{r+1}|$, so
\[
  |A_{r+1}| \equiv w - |A_r| \pmod 2.
\]
Hence $T^2$ preserves the parity $p$ of $|A|$; for even $w$ so does $T$, for odd $w$ the map $T$
exchanges the two sectors and $T^2$ has the same nonzero spectrum on both. We write
$\lambda^{(p)}_w$ for the square root of the leading eigenvalue of $T^2$ on sector $p$ of the
cylinder, and $f_w^{(p)} = \log \lambda^{(p)}_w / w$. On the cylinder $|A|$ is the number of
dimers crossing a closed loop around the axis, and for even $w$ its parity is conserved: this is
the topological $\Z_2$ sector familiar from dimer models on non-bipartite lattices
\cite{moessnersondhi2001,yaokivelson2012}. For even $w$ the two sector values differ, and they
bracket $s$ in all our data: $f^{(0)}_w > s > f^{(1)}_w$.

For free strips only the sector of $e_0$ enters $Z$, and the spectral radius of $T$ is not the
growth rate $\lambda_w$ in general: for $w = 6, 10, 14, 18$ the other sector has the larger
eigenvalue (e.g.\ $w = 6$: $\log\lambda = 2.909568954808860$ there, against
$a_6 = 2.909403518708946$). The difference of the two sectors of the free strip of even width
$k$ alternates in sign with period $4$ and shrinks by a factor $13.84$ per two columns
($k \le 18$).

\subsection{Computation}
We use float64 power iteration on $T^2$ restricted to one sector, from a random positive start
vector (fixed seed), normalised every step, stopped when the estimate of $\log \lambda$ changes
by less than $4\cdot10^{-15}$ for three consecutive iterations (19 iterations for $w \ge 20$).
The changes of the estimate shrink by a factor $0.115$ per iteration (the gap of
\S\ref{sec:z2}), so the truncation error after stopping is far below one ulp of $\log\lambda$.
Afterwards the Collatz--Wielandt bracket $\min_i (T^2 v)_i / v_i \le \lambda^2 \le \max_i (T^2
v)_i / v_i$ \cite{collatz1942,wielandt1950} over the support of $v$ is evaluated; for
$w = 20, \dots, 25$ it closes to within one ulp of $\log \lambda$. Computed without directed
rounding, it is a diagnostic here, not a certificate (the certificates are in
\S\ref{sec:float}). Widths $w \le 19$ take minutes; $w = 25$ (one sector suffices for odd $w$)
works on vectors of $2^{28}$ doubles and needs $8.4$\,GiB and $22.5$ minutes on four cores.

\subsection{No Pfaffian orientation}
An orientation of a graph is \emph{Pfaffian} if all perfect matchings have the same sign in the
Pfaffian of the oriented adjacency matrix, which then equals the number of perfect matchings up to
sign \cite[Ch.~8]{lovaszplummer1986}. Planar graphs have Pfaffian orientations
\cite{kasteleyn1961}, but so do some non-planar graphs, so the non-planarity of the king's graph
does not by itself rule out Kasteleyn's method.

\begin{theorem}\label{thm:pfaff}
Let $KL$ be even. The king board $B(K, L)$ has a Pfaffian orientation if and only if
$\min(K, L) \le 2$.
\end{theorem}
\begin{proof}
If $\min(K, L) = 1$ the board is a path. The board $B(2, L)$ is planar: each unit square with
its two diagonals is a $K_4$, consecutive squares share a vertical edge, and gluing two planar
graphs along an edge that lies on the outer face of both gives a planar graph. Both have
Pfaffian orientations \cite{kasteleyn1961}.

Let $\min(K, L) \ge 3$. A subgraph $H$ of $G$ is \emph{conformal} if $H$ and $G - V(H)$ both have
perfect matchings. A Pfaffian orientation of $G$ restricts to a Pfaffian orientation of every
conformal $H$: fix a perfect matching $R$ of $G - V(H)$; the sign of $M \cup R$ in $G$ is the sign
of $M$ in $H$ times a factor that depends only on $R$ and on the vertex order, so all perfect
matchings $M$ of $H$ have the same sign. If $L$ is even, the corner $B(3, 4)$ (rows $0, 1, 2$,
columns $0, \dots, 3$) is conformal in $B(K, L)$: the rest of the first three rows has even width
$L - 4$, the other rows have even length $L$, and horizontal dominoes cover both. If $L$ is odd,
then $K$ is even and $K \ge 4$, and the corner $B(4, 3)$ is conformal: horizontal dominoes cover
the rest of the first four rows (width $L - 3$, even) and vertical dominoes the last $K - 4$
rows. Since $B(4, 3) \cong B(3, 4)$, it remains to show that $B(3, 4)$ has no Pfaffian
orientation.

Fix the orientation $D_0$ of every edge from the lower to the higher vertex index, and a perfect
matching $M_0$. Reversing the edges of a set $X$ multiplies the sign of a perfect matching $M$ by
$(-1)^{|X \cap M|}$, so a Pfaffian orientation exists if and only if the system
\[
  \sum_{e \in M \triangle M_0} x_e \equiv \epsilon(M) + \epsilon(M_0) \pmod 2
  \qquad \text{for all perfect matchings } M
\]
has a solution $x \in \{0, 1\}^E$, where $(-1)^{\epsilon(M)}$ is the sign of $M$ under $D_0$. For
$B(3, 4)$, with $29$ edges and $43$ perfect matchings, the system is inconsistent; Gaussian
elimination over $\mathrm{GF}(2)$ finds the contradiction after $23$ of the equations.
\end{proof}

The elimination is done in exact integer arithmetic by \file{src/pfaff\_rigor.py}, which also
confirms that $B(2, L)$ is Pfaffian for $L \le 12$, that $B(3, 6)$, $B(4, 4)$, $B(4, 5)$ and
$B(4, 6)$ are not, and, for eleven boards up to $B(11, 12)$, that the matchings extending a fixed
domino tiling of the complement of the corner already give an inconsistent subsystem.

\begin{remark}
The cylinder of circumference $w \ge 3$ with $K \ge 3$ rows ($Kw$ even) and the corresponding
torus contain $B(K, w)$ as a spanning, hence conformal, subgraph, so they have no Pfaffian
orientation either. Theorem~\ref{thm:pfaff} excludes a single Pfaffian; it does not exclude
other determinantal formulas, and it does not imply hardness. The perfect matchings of a graph
of genus $g$ are a linear combination of $4^g$ Pfaffians \cite{galluccioloebl1999,tesler2000};
with $V = KL$ vertices and $E = 4KL - 3K - 3L + 2$ edges, Euler's formula gives
$g \ge (E - 3V + 6)/6 = ((K-3)(L-3) - 1)/6$, that is $g \ge 74$ for the $24 \times 24$ board. The
signed sum of \cite{wildeboerseidel2020} is a single Pfaffian.
\end{remark}

\section{The value of \texorpdfstring{$s$}{s}}\label{sec:value}

\subsection{Data}
Table~\ref{tab:odd} lists the odd cylinders. The values $f_w$ increase to $s$; the differences
$d_w = f_w - f_{w-2}$ decay by a ratio $r_w = d_{w-2}/d_w$ that drifts from $13.4$ at $w = 9$ to
$10.0$ at $w = 25$ (the table starts at $w = 13$).
The geometric (Aitken) estimate $s_w = f_w + d_w/(r_w - 1)$ settles at $0.5655119110413253$.
For even $w$ the sector values bracket $s$:
\[
\begin{array}{rll}
  w = 18: & [0.5655119060934878,\ 0.5655119160322841], \\
  w = 20: & [0.5655119102001319,\ 0.5655119118867463], \\
  w = 22: & [0.5655119108963251,\ 0.5655119111867462], \\
  w = 24: & [0.5655119110160359,\ 0.5655119110666571], & \text{midpoint } 0.5655119110413465,
\end{array}
\]
and the midpoints converge to $s$ at the same rate as the odd widths. They come from other
cylinders, but from the same code; an independent code is compared at the end of this section.

\begin{table}[ht]
\caption{Odd cylinders: $f_w = \log\lambda_w / w$, $d_w = f_w - f_{w-2}$, $r_w = d_{w-2}/d_w$ and the
geometric extrapolation $s_w$, from the float64 eigenvalues; their error, at most $6.2\cdot10^{-17}$
per site, can move the last digit of $f_w$ by one. At $w = 25$ the extended-precision inputs give
$d_{25} = 5.991\cdot10^{-14}$, $r_{25} = 9.998$ and $s_{25} = 0.56551191104132523$.}\label{tab:odd}
\begin{tabular}{rllrl}
\toprule
$w$ & $f_w$ & $d_w$ & $r_w$ & $s_w$ \\
\midrule
13 & 0.5655119032044848 & $7.943\cdot10^{-8}$  & 11.68 & 0.5655119106414456 \\
15 & 0.5655119103137916 & $7.109\cdot10^{-9}$  & 11.172 & 0.5655119110126792 \\
17 & 0.5655119109720268 & $6.582\cdot10^{-10}$ & 10.801 & 0.5655119110391899 \\
19 & 0.5655119110345892 & $6.256\cdot10^{-11}$ & 10.521 & 0.5655119110411600 \\
21 & 0.5655119110406597 & $6.071\cdot10^{-12}$ & 10.306 & 0.5655119110413121 \\
23 & 0.5655119110412587 & $5.989\cdot10^{-13}$ & 10.136 & 0.5655119110413242 \\
25 & 0.5655119110413186 & $5.994\cdot10^{-14}$ &  9.992 & 0.5655119110413253 \\
\bottomrule
\end{tabular}
\end{table}

\subsection{Extrapolation models and error bar}
All extrapolations below use the cylinders recomputed in x87 extended precision (64-bit mantissa):
the odd widths $w = 5, \dots, 25$ and both sectors of the even widths $w = 4, \dots, 24$. Float64
enters only through a separate rounding term.

\emph{Main model.} The corrections come from the gap inside a sector (\S\ref{sec:z2}): at $w = 20$
ARPACK gives $m = \log(\lambda_1/|\lambda_3|) = 1.082727$ per row, and the power iterations give the
same value at every $w = 20, \dots, 25$. For a massive correction the natural form is that of
L\"uscher \cite{luscher1986},
\[
  d_w = C e^{-mw} w^{-3/2} \Bigl(1 + \frac{c_1}{w}\Bigr),
\]
with $m$ fixed at the gap and $C, c_1$ fitted through $d_{W-2}, d_W$ (L1), or with a further term
$c_2/w^2$ through three differences (L2). The tail $s - f_W = \sum_{j \ge 1} d_{W+2j}$ is then
summed. As checks we use the generic models of the first version of this work: geometric (G,
constant ratio $r_W$); $d_w = A e^{-w/\xi} w^b$ through three points (PE) or by least squares on
six (PE6); $r_w = r_\infty + c/w$ (R1) and $r_w = r_\infty + c/w + e/w^2$ (R2); two exponentials
by Prony's method (2E); and Wynn's $\varepsilon$ algorithm \cite{wynn1956} on $f_w$, its last
full diagonal entry (WY) and the one before (WY$-$). The same models are applied to the even-width
sector midpoints, a second sequence from other cylinders computed by the same code.

\begin{table}[ht]
\caption{Estimates of $s$ by model on extended-precision inputs: from the odd widths up to
$W = 25$ and from the even midpoints up to $W = 24$, with the backtest change of each, the
difference from the estimate with the last width removed.}
\label{tab:models}
\begin{tabular}{lllll}
\toprule
model & odd, $W = 25$ & midpoints, $W = 24$ & \multicolumn{2}{l}{backtest change} \\
 & & & odd & midpoints \\
\midrule
L1   & 0.56551191104132532 & 0.56551191104132532 & $1.4\cdot10^{-17}$ & $6.5\cdot10^{-17}$ \\
L2   & 0.56551191104132532 & 0.56551191104132532 & $1.5\cdot10^{-18}$ & $3.5\cdot10^{-18}$ \\
G    & 0.56551191104132523 & 0.56551191104132564 & $1.0\cdot10^{-15}$ & $3.5\cdot10^{-15}$ \\
PE   & 0.56551191104132532 & 0.56551191104132531 & $1.7\cdot10^{-17}$ & $8.1\cdot10^{-17}$ \\
PE6  & 0.56551191104132533 & 0.56551191104132528 & $6.7\cdot10^{-17}$ & $6.0\cdot10^{-16}$ \\
R1   & 0.56551191104132532 & 0.56551191104132531 & $4.4\cdot10^{-17}$ & $1.9\cdot10^{-16}$ \\
R2   & 0.56551191104132532 & 0.56551191104132532 & $6.1\cdot10^{-19}$ & $1.6\cdot10^{-17}$ \\
2E   & 0.56551191104132531 & 0.56551191104132534 & $6.0\cdot10^{-17}$ & $2.5\cdot10^{-16}$ \\
WY   & 0.56551191104132532 & 0.56551191104132532 & $4.0\cdot10^{-18}$ & $5.4\cdot10^{-17}$ \\
WY$-$& 0.56551191104132532 & 0.56551191104132528 & $7.2\cdot10^{-18}$ & $9.3\cdot10^{-17}$ \\
\bottomrule
\end{tabular}
\end{table}

\emph{Predictive backtest.} A model is kept if its backtest change is below $1\%$ of the last
correction $|d_W|$ of its sequence ($6.0\cdot10^{-16}$ for the odd widths, $1.9\cdot10^{-15}$ for
the midpoints). Every model passes on both sequences except G, which ignores the drift of $r_w$
(the PE fit attributes it to a power prefactor, $\xi_{\mathrm{odd}} = 0.929$, $b = -1.80$). Measured
against $|d_W|$, the backtest changes of L1 stay at $2$--$4\cdot10^{-5}$ on the odd widths for every
cut $W = 19, 21, 23$, and those of L2 and R2 fall to $1$--$2\cdot10^{-6}$ at $W = 23$. The $18$ kept estimates lie in
$[0.56551191104132527,\ 0.56551191104132534]$.

\emph{Error budget.} We add four terms:
\begin{enumerate}
\item model form, $9.0\cdot10^{-17}$: the half-spread of the kept estimates ($3.0\cdot10^{-17}$) plus
the largest predicted remaining error, each model's backtest change scaled by the ratio
$|d_W/d_{W-2}|$ of the last two corrections ($6.0\cdot10^{-17}$);
\item a hidden slow mode, $3.4\cdot10^{-17}$: we add a term $B x^{(w-25)/2}$ with
$x = e^{-2/1.2351}$, the slowest decay seen anywhere in the data (the parity splitting of
\S\ref{sec:z2}), and fit it together with L1 ($C$, $c_1$ and $B$) by relative least squares on
$w = w_0, \dots, 25$ for $w_0 = 11, 13, 15, 17$. The fitted $B$ depends on the window, between
$0.9\cdot10^{-3}\,d_{25}$ ($w_0 = 11$) and $1.95\cdot10^{-3}\,d_{25}$ ($w_0 = 13$). We allow every $B$ within three
standard errors of a fitted value, over all windows and widened to include $B = 0$:
$-1.6\cdot10^{-3}\,d_{25} \le B \le 3.4\cdot10^{-3}\,d_{25}$. Since $s$ is linear in $B$, the larger shift at
the two ends is the term. The next odd width does not show such a mode: the independent code of
\file{xcheck/}, described at the end of this section, gives $d_{27}$ within $1.5\cdot10^{-4}$ of the L1 prediction from $W = 25$, where a mode
with $B = 1.95\cdot10^{-3}\,d_{25}$ would move it by $1.0\cdot10^{-3}$. This is consistent with the fitted $B$
taking up the misfit of L1 at small $w$, though it rests on one width. Read as a mode, the whole
deviation allows $|B| \le 2.8\cdot10^{-4}\,d_{25}$, which would move the main estimate by at most
$1.4\cdot10^{-18}$; this check is not part of the budget;
\item the gap, $6\cdot10^{-20}$: the change of L1 when $m$ runs over the measured gaps;
\item rounding, $1.5\cdot10^{-16}$: the largest change of a kept model between float64 and
extended-precision inputs (the float64 values of $\log\lambda_w$ differ by at most $6.2\cdot10^{-17}$
per site over $33$ cylinders).
\end{enumerate}
The total is $2.7\cdot10^{-16}$. We quote a little more than twice this,
\[
  \boxed{\,s = 0.5655119110413253 \pm 6\cdot10^{-16}\,}
\]
with the main model L1 on the odd widths, $s = 0.5655119110413253176$, as the centre. This is an
estimate, not a bound: the budget assumes that no correction slower than those in the data
appears beyond $w = 25$.

\emph{Further digits.} Each further digit of the truncation error needs about $\log 10/m = 2.13$
more widths. The state vector doubles with each width, and with it the peak memory of
\file{src/cylw.py} ($4.3$\,GiB at $w = 24$, $8.4$\,GiB at $w = 25$), so a digit costs about
$2^{2.13} = 4.4$ times the memory, and the run time grows by about the same factor. The quoted error
bar would not follow: its largest term is the rounding term ($1.5\cdot10^{-16}$ of
$2.7\cdot10^{-16}$), which more widths do not reduce, and the extended-precision inputs agree with
the independent code below to $7.8\cdot10^{-20}$ per site, so about $19$ digits is the end of the
$64$-bit mantissa of the x87 \texttt{long double}. Restricting the full row transfer to zero momentum and to states even under reflection would
divide the state space by about $2w$, but the cell steps used here break the rotation symmetry of the
intermediate states, so this needs a different implementation. Beyond a few more digits the natural
tools are the corner transfer matrices and their renormalization-group variant mentioned in the
introduction, with the parity sector built into the encoding.

\emph{Cross-checks of the float64 eigenvalues.} Exact integer ratios
$Z_{\mathrm{cyl}}(K+2, w)/Z_{\mathrm{cyl}}(K, w)$ for $w = 9, 11, 12, 13$ agree with the power
iteration to $4.1\cdot10^{-16}$ in $\log \lambda$, and ARPACK on the full $w = 20$ transfer
reproduces both sector values to $1.8\cdot10^{-15}$. The exact counts of the square boards give an
independent, eigenvalue-free check of $s$, discussed with the corner free energy in
\S\ref{sec:corner}.

\emph{Independent recomputation.} A transfer-matrix code written from scratch to check this paper
(\file{xcheck/}; nothing in it is taken from \file{src/}) reproduces the cylinder values per site to
$7.8\cdot10^{-20}$ against the extended-precision inputs and to $6.2\cdot10^{-17}$ against the float64
ones. It also goes beyond them: both sectors of $w = 26$ and the odd width $w = 27$ in extended
precision, and both sectors of $w = 28$ in float64. The new differences $d_{26}$ and $d_{27}$ agree
with the L1 predictions made at $W = 24$ (midpoints) and $W = 25$ (odd widths) to $1.9\cdot10^{-4}$ and
$1.5\cdot10^{-4}$ relative, and $d_{28}$, whose float64 inputs resolve it only to about
$6\cdot10^{-3}$, to $5.4\cdot10^{-4}$; L1 on the odd widths up to
$W = 27$ gives $0.5655119110413253187$, $1.1\cdot10^{-18}$ from the centre above. These widths test the
extrapolation outside its sample; they are not used in the estimate.

\section{Rigorous bounds}\label{sec:bounds}

\begin{theorem}[computer-assisted]\label{thm:bounds}
$0.545343749727 \le \underline s \le \bar s \le 0.582767477484$.
\end{theorem}

For comparison, the Kahn--Lov\'asz bound (the number of perfect matchings of a graph is at most
$\prod_v (d_v!)^{1/(2 d_v)}$; unpublished by its authors, see \cite{cucklerkahn2009}, with a
proof in \cite{alonfriedland2008}), with all degrees $d_v \le 8$, gives
$\bar s \le \log(8!)/16 = 0.662787681422$.

\emph{Trust base.} The proof of Theorem~\ref{thm:bounds} is computer-assisted, and its correctness
rests on two things beyond the mathematics below. First, the certificates are evaluated in
IEEE-754 binary64 arithmetic with round-to-nearest (no fused multiply-add contraction, no flush to
zero), and their validity follows from the rounding analysis of \S\ref{sec:float}. Second, the
transfer matrix must be the right matrix: the code is argued correct and validated against
brute-force enumeration on small boards (\S\ref{sec:float}, \S\ref{sec:checks}), but it has not been
formally verified.

Both bounds of Theorem~\ref{thm:bounds} come from strips of width $w = 24$; the upper bound uses
the run-weighted transfer of Lemma~\ref{lem:runs}. Table~\ref{tab:bounds} shows how the bounds
improve with $w$. Since $\log \lambda_w / w = s + g(w)/w$ with $g(w) \to -2\sigma$
(\S\ref{sec:surface}), the lower bound approaches $s$ like $2\sigma/w = 0.484/w$; the upper bound
with the uniform weight $\sqrt\varphi$ approaches it, empirically, like $0.497/w$, and the
run-weighted one like $0.414/w$.

\begin{table}[ht]
\caption{Certified bounds from strips of width $w$: the lower bound, the upper bound with the
uniform side weight $\sqrt\varphi$ (Lemma~\ref{lem:strip}) and the run-weighted upper bound
(Lemma~\ref{lem:runs}).}\label{tab:bounds}
\begin{tabular}{rlll}
\toprule
$w$ & lower bound on $\underline s$ & upper, weight $\sqrt\varphi$ & upper, run weights \\
\midrule
4  & 0.445893077494 & 0.689765394232 & 0.669090528409 \\
8  & 0.505014525148 & 0.627626147026 & 0.617278850986 \\
12 & 0.525175645897 & 0.606921301094 & 0.600023045982 \\
16 & 0.535259669620 & 0.596568952598 & 0.591395260725 \\
20 & 0.541310117470 & 0.590357544276 & 0.586218590772 \\
24 & 0.545343749727 & 0.586216605404 & 0.582767477484 \\
\bottomrule
\end{tabular}
\end{table}

\subsection{Lower bound}

\begin{lemma}[superadditivity]\label{lem:super}
For even $K$ and $L = kw + r$ with $0 \le r < w$: $Z(K, L) \ge Z(K, w)^k$.
\end{lemma}
\begin{proof}
Perfect matchings of $k$ strips $B(K, w)$ placed side by side, completed by vertical dominoes in
the last $r$ columns ($K$ is even), are distinct perfect matchings of $B(K, L)$.
\end{proof}

\begin{lemma}[Collatz--Wielandt, lower]\label{lem:cwlow}
Let $A = T^2$ be the free-strip transfer of width $w$, so that $Z(2j, w) = (A^j e_0)_0$. Let
$v \ge 0$ with $v_0 > 0$, $\supp v \subseteq \supp A^n e_0$ for some $n$, and $Av \ge c\,v$
componentwise. Then $Z(2(n + j), w) \ge c^j v_0 / C'$ for all $j$, where
$C' = \max_{i \in \supp v} v_i / (A^n e_0)_i < \infty$. In particular $\lambda_w^2 \ge c$.
\end{lemma}
\begin{proof}
$v \le C' A^n e_0$ componentwise and $A \ge 0$, so
$Z(2(n+j), w) = (A^j A^n e_0)_0 \ge (A^j v)_0 / C' \ge c^j v_0 / C'$.
\end{proof}

\begin{proof}[Proof of the lower bound]
By Lemma~\ref{lem:super}, for even $K$,
\[
  \frac{\log Z(K, L)}{KL} \ge \frac{\lfloor L/w\rfloor}{L}\,\frac{\log Z(K, w)}{K}
  \longrightarrow \frac{\log\lambda_w}{w} \qquad (K, L \to \infty);
\]
for odd $K$ (then $L$ is even) exchange the roles of the sides. So $\underline s \ge a_w / w$ for every $w$, and
$\underline s \ge s^*$. The certificate of \S\ref{sec:float} gives
$a_{24}/24 \ge \frac{1}{48}\log c_{24} \ge 0.545343749727$.
\end{proof}

\subsection{Upper bound}
Cut $B(K, kw)$ into $k$ strips $S^{(1)}, \dots, S^{(k)}$ of width $w \ge 2$ (columns
$jw, \dots, jw + w - 1$). King edges change the column by at most one, so every edge of a perfect
matching $X$ lies inside a strip or joins column $jw - 1$ to column $jw$ (an edge across cut $j$).
Let $A_j$ be the rows of the cells of column $jw - 1$ matched across cut $j$ (right side of
$S^{(j)}$), $B_{j+1}$ the rows of the cells of column $jw$ matched across cut $j$ (left side of
$S^{(j+1)}$), and $B_1 = A_k = \emptyset$. Let $N(B, A)$ be the number of perfect matchings of
\[
  B(K, w) \setminus \bigl(\{(b, 0) : b \in B\} \cup \{(a, w-1) : a \in A\}\bigr),
\]
and $M(A, B)$ the number of perfect matchings of the bipartite graph on $A \sqcup B$ with
$a \sim b$ iff $|a - b| \le 1$.

\begin{lemma}[cut identity]\label{lem:cut}
$Z(K, kw) = \sum \prod_{j=1}^{k} N(B_j, A_j) \prod_{j=1}^{k-1} M(A_j, B_{j+1})$, the sum over all
$A_1, B_2, A_2, \dots, B_k$ with $B_1 = A_k = \emptyset$.
\end{lemma}
\begin{proof}
Since $w \ge 2$ the sets are well defined. The restriction of $X$ to strip $j$ is a perfect matching
of $S^{(j)} \setminus (B_j \cup A_j)$, and its restriction to cut $j$ a perfect matching between
$A_j$ and $B_{j+1}$; conversely every such tuple of pieces is a perfect matching of $B(K, kw)$.
\end{proof}

The same proof applies to strips of unequal widths, and to a last strip of width $1$, whose cells
can only be matched across its left cut (there is no cut to its right); the proof of the upper
bound uses this for the remainder strip.

\begin{lemma}[crossing matchings]\label{lem:fib}
If $|A| = |B| = n$ then $M(A, B) \le F_{n+1} \le \varphi^n$, where $F_1 = F_2 = 1$ are the
Fibonacci numbers and $\varphi = (1 + \sqrt 5)/2$; $M(A, B) = 0$ if $|A| \ne |B|$.
\end{lemma}
\begin{proof}
Let $a_1 < \dots < a_n$, $b_1 < \dots < b_n$ and $M(n) = \max M(A, B)$. \emph{$a_1$ is matched to
$b_1$ or $b_2$}: if $a_1 \sim b_j$ with $j \ge 3$, then $b_1 \le b_3 - 2 \le a_1 - 1$, while every
$a \ne a_1$ has $a \ge a_1 + 1 \ge b_1 + 2$, so $b_1$ has no partner. \emph{If $a_1 \sim b_2$},
then $b_1 \le b_2 - 1 \le a_1$ needs a partner $a_i \ge a_1 + 1$ with $a_i \le b_1 + 1 \le a_1 + 1$,
so $a_i = a_1 + 1$, $b_1 = a_1$, and the pair $a_2 b_1$ is forced. Thus every matching contains
$a_1 b_1$ (the rest is a matching of $n - 1$ against $n - 1$ cells) or both $a_1 b_2$ and
$a_2 b_1$ (the rest is $n-2$ against $n-2$), so $M(n) \le M(n-1) + M(n-2)$ with
$M(0) = M(1) = 1$, i.e.\ $M(n) \le F_{n+1}$. Finally $F_{n+1} \le \varphi^n$ by induction.
\end{proof}

Equality holds when $A = B$ is an interval. The same branching gives more when $A$ is not an
interval. A \emph{run} of a set $S$ of rows is a maximal block of consecutive rows in $S$; put
$F(S) = \prod_{\text{runs } R \text{ of } S} F_{|R|+1}$.

\begin{lemma}[crossing matchings, run form]\label{lem:runs}
$M(A, B) \le \min(F(A), F(B))$, with equality $M(S, S) = F(S)$. Moreover $F_2 = 1$ and
$F_{l+1} \le 2\varphi^{l-2}$ for $l \ge 2$.
\end{lemma}
\begin{proof}
The relation $a \sim b$ is symmetric, so it suffices to show $M(A, B) \le F(A)$, by induction on
$n = |A| = |B|$ (the case $n = 0$ is trivial, and $M = 0$ if $|A| \ne |B|$). Let $a_1$ be the
smallest row of $A$ and $l$ the length of its run. By the proof of Lemma~\ref{lem:fib}, every
matching contains $a_1 b_1$, or contains $a_1 b_2$ and $a_2 b_1$ with $a_2 = a_1 + 1 \in A$; the
second case needs $l \ge 2$. Removing $a_1$ shortens the first run to $l - 1$, removing $a_1, a_2$
shortens it to $l - 2$ (a run of length $0$ disappears, and $F_1 = 1$), and the other runs are
unchanged; write $R$ for the product of their factors. By induction
$M(A, B) \le F_l R + F_{l-1} R = F(A)$ if $l \ge 2$, and $M(A, B) \le F(A \setminus \{a_1\}) = F(A)$
if $l = 1$. For $A = B = S$, two runs of $S$ are at distance at least $2$, so no edge joins
different runs, and on a run of length $l$ the count is $F_{l+1}$ (Lemma~\ref{lem:fib}). Finally
$F_3 = 2$, $F_4 = 3 \le 2\varphi$ and $F_{l+2} = F_{l+1} + F_l \le 2\varphi^{l-3}(\varphi + 1) =
2\varphi^{l-1}$.
\end{proof}

Hence $M(A, B) \le \sqrt{F(A) F(B)} \le h(A)\, h(B)$ for every weight $h$ with $h(S)^2 \ge F(S)$.
We use weights that are products over the runs of $S$ of per-cell factors depending on the position
in the run: $1, \sqrt2, \sqrt\varphi, \sqrt\varphi, \dots$, so that a run of length $l$ gets
$h^2 = 1, 2, 2\varphi, 2\varphi^2, \dots$, which is at least $F_{l+1}$ by Lemma~\ref{lem:runs}. The
uniform weight $h(S) = \varphi^{|S|/2}$ of Lemma~\ref{lem:fib} is the special case of one factor
$\sqrt\varphi$ per cell; it overpays on short runs, and for $A = B = \{0, 2, 4, \dots\}$, a set of
runs of length $1$, it gives $\varphi^{|S|}$ where $F(S) = 1$. The smallest admissible weight is
$h = \sqrt F$ itself, with factors $\sqrt{F_{t+2}/F_{t+1}}$ for the $t$-th cell of a run; it needs
the full run length as a state, and the three factors above truncate it after the second cell.

Put $z = \sqrt\varphi$ and
\[
  W(K, w) = \sum_{A, B} N(B, A)\, z^{|A| + |B|},
\]
the weighted number of matchings of a strip in which a cell of the first or last column may be
``covered from outside'' at weight $z$.

\begin{lemma}[strip inequality]\label{lem:strip}
$Z(K, kw) \le W(K, w)^k$. More generally $Z(K, kw) \le W_h(K, w)^k$ with
$W_h(K, w) = \sum_{A, B} N(B, A)\, h(A)\, h(B)$, for every $h$ with $h(S)^2 \ge F(S)$.
\end{lemma}
\begin{proof}
In Lemma~\ref{lem:cut} replace $M(A_j, B_{j+1})$ by $\varphi^n = z^{|A_j|} z^{|B_{j+1}|}$
(Lemma~\ref{lem:fib}). The factor $z^{|A_j|}$ belongs to strip $j$ and $z^{|B_{j+1}|}$ to strip
$j+1$, so each summand is at most $\prod_j N(B_j, A_j) z^{|B_j| + |A_j|}$. Dropping the constraints
$|A_j| = |B_{j+1}|$ and $B_1 = A_k = \emptyset$ adds only nonnegative terms, and the sum
factorises over the strips. For general $h$ replace $\varphi^n$ by $h(A_j)\, h(B_{j+1})$
(Lemma~\ref{lem:runs}); $h(\emptyset) = 1$.
\end{proof}

The transfer of \S\ref{sec:model} computes $W$ when a cell in column $0$ or $w-1$ may also be
covered from outside at weight $z$: $W(K, w) = e_0^\top T_z^K e_0$, $T_z \ge 0$. For the run weights
the transfer carries, for each of the two side columns, the position of the previous row's side
cell in a run of outside covers ($0$: not covered from outside, $1$: first cell of a run, $2$: a
later cell). A side cell covered from outside gets the factor of its position, $1$, $\sqrt2$ or
$\sqrt\varphi$, and advances the state; any other move of a side cell resets it to $0$. The state
space grows by a factor $9$, and $W_h(K, w) = \sum_{l, r} (T_h^K e_{0})_{(l, r, 0)}$ with
$T_h \ge 0$. Lemma~\ref{lem:cwup} holds for $T_h$ with the same proof, applied to each of the nine
final entries: $W_h(2j, w) \le c' C^j$ with $c' = \sum_{l, r} v_{(l, r, 0)} / v_0$, a constant
independent of $j$.

\begin{lemma}[Collatz--Wielandt, upper]\label{lem:cwup}
If $A = T_z^2$, $v \ge 0$, $v_0 > 0$ and $Av \le Cv$ componentwise (including $(Av)_i = 0$
where $v_i = 0$), then $W(2j, w) \le C^j$ for all $j$.
\end{lemma}
\begin{proof}
$e_0 \le v/v_0$ and $A \ge 0$, so $A^j e_0 \le A^j v / v_0 \le C^j v / v_0$; take the entry $0$.
\end{proof}

\begin{proof}[Proof of the upper bound]
Let $K, L \to \infty$ with $KL$ even. If $K$ is odd then $L$ is even and $Z(K, L) \le Z(K+1, L)$
(add a row of horizontal dominoes), so
$(KL)^{-1}\log Z(K, L) \le \frac{K+1}{K}\,((K+1)L)^{-1}\log Z(K+1, L)$ and the factor
$(K+1)/K$ tends to $1$; we may therefore take $K = 2j$. Write $L = kw + r$, $0 \le r < w$.
Cutting off a last strip of width $r$ as in Lemma~\ref{lem:strip} gives
$Z(K, L) \le W_h(K, w)^k W_r(K)$ ($h$ the uniform or the run weight), where $W_r$ counts the configurations of a width-$r$ strip with
weighted outside covers; crudely $W_r(K) \le (5 + 2\sqrt2)^{Kr}$ (per cell: already covered, one of
at most four later partners, or one of at most two outside covers, each of weight at most $\sqrt2$).
By Lemma~\ref{lem:cwup}, with the constant $c' \ge 1$ above for the run weights ($c' = 1$ for
$\sqrt\varphi$),
\[
  \frac{\log Z(K, L)}{KL} \le \frac{k}{L}\,\frac{\log C_w}{2} + \frac{k \log c'}{KL}
  + \frac{w \log(5 + 2\sqrt2)}{L} \longrightarrow \frac{\log C_w}{2w}.
\]
The certificate of \S\ref{sec:float} at $w = 24$ gives $\bar s \le 0.586216605404$ with the weight
$\sqrt\varphi$ and $\bar s \le 0.582767477484$ with the run weights.
\end{proof}

\subsection{Floating-point certificates}\label{sec:float}
The certificates use float64 vectors $v$ produced by power iteration; how $v$ was obtained does not
matter. For the lower bound $v$ is the normalised iterate $A^n e_0$; floating point never creates a
positive entry from an exact zero (all terms are nonnegative), so $\supp v \subseteq \supp A^n e_0$.
The code computes $y = \fl(Av)$, asserts $v_0 > 0$, $v_i > 10^{-280}$ on $\supp v$ and (upper
bound) $y_i = 0$ off $\supp v$, and sets
\[
  c = \min_{i \in \supp v} \fl(y_i / v_i)\,(1 - 10^{-12}), \qquad
  C = \max_{i \in \supp v} \fl(y_i / v_i)\,(1 + 10^{-12}).
\]
Each cell step writes every entry as a sum of at most six nonnegative terms (already covered,
outside cover times $z$, right, down-left, down, down-right), so over the $2w$ steps of $A$,
$\fl(Av)_i \in (Av)_i\,[(1-u)^{12w}, (1+u)^{12w}]$ with $u = 2^{-53}$: a relative error of at most
$12wu\,(1 + o(1)) = 3.2\cdot10^{-14}$ at $w = 24$, plus one rounding in the division. All weights
are at least $1$ ($z$ is rounded up, so $T_{\fl(z)} \ge T_{\sqrt\varphi}$ entrywise) and the
entries on the support exceed $10^{-280}$, so there is no underflow and a float entry vanishes iff
the exact one does. Since $10^{-12} \gg 3.2\cdot10^{-14} + u$, the inequalities $Av \ge cv$ and
$Av \le Cv$ hold exactly. For the run-weighted transfer with $n = 3$ states per side, an interior
cell step is unchanged; at the left side column an entry of run state $0$ collects at most $4n$ terms, and at the
right one at most $3n$, while an entry of run state $1$ or $2$ is reached only by outside covers and
needs at most two products by a factor and one addition. Over the $2w$ steps of $A$ this is at most
$2\bigl(4(w-2) + (4n-1) + (3n-1)\bigr) = 8w + 14n - 20 \le 12w$ roundings for $w \ge 6$, the same
budget. The factors are rounded up to floats, so the float transfer dominates the exact one
entrywise, and the factor $1$ is exact. The logarithms are taken with 30-digit arithmetic and the results are
rounded outward to 12 decimals.

\emph{Audit.} The following were checked in exact arithmetic: the cut identity and both strip
inequalities on $58$ boards ($K \le 8$, $kw \le 12$); $\max M(A, B) = F_{n+1}$ for all $A, B$ in
columns of length $m \le 11$, and the branching step for $m \le 8$; the transfer with weight $z$
against brute force as a polynomial in $z$ ($162$ evaluations); and, for $w = 2, \dots, 12$, the
float ratios against exact rational ratios for the same float vector ($|\text{float}/\text{exact}
- 1| \le 1.7\cdot10^{-16}$). The audit is sensitive: with $z = 0$ the strip inequality fails on all
$40$ boards with $k \ge 2$. On these small boards $z = 0.5$ would suffice. The proof needs only
$h(S)^2 \ge F(S)$ for every $S$, which $\sqrt\varphi$ meets uniformly and the run weights meet with
less to spare. For the run weights, also exactly: $M(A, B) \le
\min(F(A), F(B))$ and $M(S, S) = F(S)$ for all $A, B$ in columns of length $m \le 8$; the strip
inequality with the run weights in interval arithmetic on the same $58$ boards; the run-weighted
transfer against brute force at rational factors ($57$ evaluations); and the float ratios against
exact ones for $w = 2, \dots, 8$ ($|\text{float}/\text{exact} - 1| \le 1.7\cdot10^{-16}$).

\section{Surface free energy}\label{sec:surface}
For the free strip of width $k$, $\log Z(K, k) = K a_k + O(1)$, and
\[
  g(k) = a_k - s k \longrightarrow -2\sigma,
\]
where $\sigma$ is the free energy per unit length of a free edge. Even $k$ approach the limit from
above, odd $k$ from below, and every even value lies above every odd one (Table~\ref{tab:surface};
$a_k$ from extended-precision power iteration for $k = 14, \dots, 24$ and float64 with a
Collatz--Wielandt bracket for $k = 25, 26$).

\begin{table}[ht]
\caption{$g(k) = \log\lambda_k - sk$ for free strips.}\label{tab:surface}
\begin{tabular}{rlrl}
\toprule
$k$ & $g(k)$ & $k$ & $g(k)$ \\
\midrule
19 & $-0.4840358718678495$ & 20 & $-0.4840358714180326$ \\
21 & $-0.4840358715709653$ & 22 & $-0.4840358715195378$ \\
23 & $-0.4840358715369125$ & 24 & $-0.4840358715310018$ \\
25 & $-0.4840358715330068$ & 26 & $-0.4840358715323304$ \\
\bottomrule
\end{tabular}
\end{table}

The odd sequence converges cleanly (all models of \S\ref{sec:value} agree to $2\cdot10^{-16}$): its
differences $g(k+2) - g(k)$ shrink by a factor $8.719$ per two columns ($k = 19, \dots, 25$), that
is $e^{2 \cdot 1.0827}$, the in-sector gap of \S\ref{sec:z2}. The even one carries two corrections,
this mode and the parity-splitting mode of the free strip, which alternates with period $4$ and
shrinks by a factor $13.84$ per two columns (\S\ref{sec:model}); and only the two-exponential and Wynn models are stable under the
backtest. The stable models give $-2\sigma \in [-0.4840358715325015, -0.4840358715325006]$; the
error of $s$, multiplied by $k = 26$, dominates:
\[
  -2\sigma = -0.484035871532501 \pm 1.8\cdot10^{-14}, \qquad
  \sigma = 0.242017935766251 \pm 9\cdot10^{-15}.
\]
The errors of $\sigma$ and $s$ are correlated: an error $\delta s$ in $s$ moves the extrapolated
$-2\sigma$ by $-26\,\delta s$, so $\delta\sigma \approx 13\,\delta s$, and quantities that depend on
both add their errors linearly.

\subsection{Corner free energy and exact counts}\label{sec:corner}
The exact counts $Z(K, L)$ of the free boards for $K, L \le 24$ (computed modulo ten primes and
reconstructed by the Chinese remainder theorem, \file{src/zexact.py}) test $s$ and $\sigma$
without eigenvalues and give the next term of the expansion,
\[
  \log Z(K, L) = sKL - 2\sigma(K + L) + \kappa + o(1).
\]
On the squares, $\rho(L) = \log Z(L, L) - sL^2 + 4\sigma L$ converges to $\kappa$ with
differences shrinking by a ratio that drifts from $5.85$ to $5.58$ over $L = 14, \dots, 24$. A
geometric (Aitken) extrapolation at $L = 24$ gives $0.363196256048$; it changes by
$1.4\cdot10^{-10}$ from $L = 22$, and by $2.6\cdot10^{-11}$ if the ratio is extrapolated linearly in
$1/L$. The errors of $s$ and $\sigma$ enter as $L^2\,\delta s + 4L\,\delta\sigma = 1.2\cdot10^{-12}$,
added linearly because they are correlated. Taking twice the larger extrapolation difference,
\[
  \kappa = 0.36319625605 \pm 3\cdot10^{-10}.
\]
Rectangles agree: $|\rho(K, L) - \kappa| \le 6.3\cdot10^{-8}$ for $\min(K, L) \ge 20$, with the
sign alternating with the parity of $\min(K, L)$.

Since A239273 counts the squares, $a(n) = Z(2n, 2n)$,
\begin{gather*}
  a(n) \sim C\, d^{\,n^2} f^{-n},\\
  C = e^\kappa = 1.437918032, \quad d = e^{4s} = 9.60273202731337, \quad f = e^{8\sigma} = 6.931966132248,
\end{gather*}
with relative error $-2.0\cdot10^{-9}$ at $n = 12$ and about a factor $5.6$ smaller per step in $n$.
The b-file of A239273 up to $n = 12$, contributed by the author, agrees with these counts term by term.
The form, a constant times one exponential in the side and one in the area, is the one McKay found
numerically for the hard squares, the independent sets of the $n \times n$ grid \cite{oeisA006506}.
The ratio of the extrapolation differences is still drifting at $L = 24$, so the slowest edge and
corner correction is not yet settled; the error bar of $\kappa$ allows for this by doubling.

The counts alone also give $s$. The second differences
$[\log Z(L+2, L+2) - 2\log Z(L, L) + \log Z(L-2, L-2)]/8$ approach $s$ from above, $5.3\cdot10^{-9}$
above it at $L = 22$, and their geometric extrapolation is $1.0\cdot10^{-10}$ above the value of
\S\ref{sec:value}. The mixed differences
$[\log Z(L+2, L+2) - \log Z(L+2, L) - \log Z(L, L+2) + \log Z(L, L)]/4$, in which the edge and
corner terms cancel exactly, approach $s$ from below ($1.1\cdot10^{-11}$ at $L = 22$), and their
extrapolation agrees with it to $5.9\cdot10^{-13}$. The square counts of the JavaScript solvers of
the repository, an independent code base, agree with these counts where both exist ($L \le 20$).

\section{Parity splitting and its decay length}\label{sec:z2}
For even $w$ let $\Delta_w = \log\lambda^{(0)}_w - \log\lambda^{(1)}_w$ and
$\xi_w = 2/\log(\Delta_{w-2}/\Delta_w)$ (Table~\ref{tab:z2}). A pure exponential
$\log\Delta_w = a - w/\xi$ fits $w = 10, \dots, 24$ poorly ($\xi = 1.1931$, rms $1.5\cdot10^{-2}$);
with a power prefactor, $\log\Delta_w = a - w/\xi + b\log w$ gives $\xi = 1.2323$, $b = -0.433$,
rms $9.3\cdot10^{-4}$, and extrapolating $1/\xi_w$ linearly in $1/w$ gives $\xi_\infty = 1.2351$
(prefactor exponent $b = -0.467$). With $b$ fixed at $-1/2$, $\xi = 1.2387$ (rms $2.6\cdot10^{-3}$).
The three estimates span $\xi \in [1.2323, 1.2387]$; we quote $\xi = 1.235 \pm 0.004$.

\begin{table}[ht]
\caption{Splitting of the parity sectors on even cylinders.}\label{tab:z2}
\begin{tabular}{rll@{\qquad}rll}
\toprule
$w$ & $\Delta_w$ & $\xi_w$ & $w$ & $\Delta_w$ & $\xi_w$ \\
\midrule
10 & $1.521094\cdot10^{-4}$ & 1.15438 & 18 & $1.788983\cdot10^{-7}$ & 1.19451 \\
12 & $2.782917\cdot10^{-5}$ & 1.17750 & 20 & $3.373229\cdot10^{-8}$ & 1.19878 \\
14 & $5.124578\cdot10^{-6}$ & 1.18201 & 22 & $6.389264\cdot10^{-9}$ & 1.20204 \\
16 & $9.544544\cdot10^{-7}$ & 1.19001 & 24 & $1.214909\cdot10^{-9}$ & 1.20486 \\
\bottomrule
\end{tabular}
\end{table}

The length $\xi$ describes the splitting of the two topological sectors, a property of the
cylinder; it is not the bulk correlation length. Inside a sector the transfer has a gap: at
$w = 20$ ARPACK gives $\log(\lambda_1/|\lambda_3|) = 1.082727$ per row, with a negative subleading
eigenvalue $\lambda_3$ (\file{out/out\_crosscheck.txt}), and the convergence of the power iteration
gives the same gap at every $w = 20, \dots, 25$. The bulk correlation length along the cylinder is
therefore $1/1.0827 = 0.92$, and the odd-width corrections, which behave like
$e^{-w/\xi_{\mathrm{odd}}} w^b$ with $1/\xi_{\mathrm{odd}} \approx 1.07$, are governed by the same
gap. The prefactor form $w^b$ of $\Delta_w$ follows the analysis of the topological splitting in
\cite{ioselevich2002, yaokivelson2012}; we take only the form from there, not their value of
$\xi$, which belongs to other models. All finite-size corrections we see are exponentially small,
as expected for a non-critical (gapped) dimer liquid. On the square lattice, which is critical, the
corrections are algebraic, $(f_w - G/\pi)w^2 \to \pi/6$, and a polynomial fit in $1/w^2$ to the
transfer-matrix values for $w \le 18$ gives $G/\pi$ only to $6\cdot10^{-8}$; the exact free-fermion
values give the same digits, so the truncated asymptotic series, not the transfer matrix, limits
this check (\file{out/out\_validate.txt}).

\section{Checks on solvable lattices}\label{sec:checks}
The transfer code has a switch for the square lattice (no diagonals) and the triangular lattice
(one diagonal). Against brute-force enumeration of perfect matchings it agrees on all $90$ boards
tested (king, triangular and square; free strips and cylinders; $w = 3, \dots, 6$, $h = 1, \dots, 4$).
It reproduces the exact domicule counts of the repository for all square boards up to
$18\times18$, for rectangles, and the $20\times20$ count.

\emph{Square lattice.} The free-strip growth rates agree with the product formula
$\lambda^{\mathrm{sq}}_w = \prod_{j=1}^{\lfloor w/2 \rfloor} \bigl(a_j + \sqrt{1 + a_j^2}\bigr)$,
$a_j = \cos(\pi j/(w+1))$ \cite{kasteleyn1961}, and both cylinder sectors with the free-fermion
formula, to $10^{-16}$.

\emph{Triangular lattice.} The triangular lattice is planar, and we found a Pfaffian orientation
by a sign pattern periodic modulo $2$, checked against the exact counts on small boards and by
$|\mathrm{Pf}| = Z$ on the $8\times8$, $10\times10$ and $12\times12$ boards ($1848358752$,
$1852300296658080$, $57261989986381118073728$). Averaging the logarithm of the Bloch determinant
over a $32\times32$ and a $64\times64$ grid of momenta gives the same value,
\[
  s_{\mathrm{tri}} = 0.428594537464959,
\]
and the transfer matrix on the cylinder of circumference $17$ gives $0.428594537464956$, a
difference of $2.9\cdot10^{-15}$. This agrees with the literature value $0.857189074\ldots$ per
dimer \cite{wu2006}; the orientation goes back to Montroll \cite{montroll1964}, and the check only
tests our code.

\emph{Context.} Table~\ref{tab:context} places $s$ among the three lattices. The Bethe value
$\frac12\log\bigl((d-1)^{d-1}/d^{d-2}\bigr)$ per vertex is the growth rate of the expected number of
perfect matchings of a random $d$-regular graph \cite{bollobasmckay1986} and the zeroth order of
Nagle's expansion \cite{nagle1966}. For $d$-regular bipartite graphs it is a lower bound
\cite{schrijver1998}, and the square lattice lies above it; on the non-bipartite triangular and
king lattices it overestimates $s$, by $2.7\%$ and $1.2\%$.

\begin{table}[ht]
\caption{Entropy per site of the three lattices, the Bethe value and the Kahn--Lov\'asz bound
$\log(d!)/(2d)$.}\label{tab:context}
\begin{tabular}{lrllll}
\toprule
lattice & $d$ & $s$ & Bethe & $s/\text{Bethe}$ & Kahn--Lov\'asz \\
\midrule
square     & 4 & $G/\pi = 0.291560904031$ & 0.261624072 & 1.1144  & 0.397256729 \\
triangular & 6 & 0.428594537465            & 0.440075843 & 0.97391 & 0.548270934 \\
king       & 8 & 0.565511911041            & 0.572360897 & 0.98803 & 0.662787681 \\
\bottomrule
\end{tabular}
\end{table}

\section{Convexity conjectures and a conditional two-sided bound}\label{sec:conj}

The bounds of Theorem~\ref{thm:bounds} converge like $1/w$ because every strip has two free
sides. Differences of strip growth rates cancel the surface term. On the square lattice,
reflection positivity turns them into rigorous bounds that converge exponentially
\cite{calkinwilf1998}; here it fails. The gluing form between two half-boards is the number
$M(A, B)$ of matchings across the cut (Lemma~\ref{lem:fib}), and it is indefinite: for columns of
length $3$ its smallest eigenvalue is $1 - \sqrt2$, on a vector realised by actual half-boards.
Nor do the cylinders give a proven upper bound. At every odd circumference $w = 3, 5, \dots, 25$ the
cylinder value $f_w$ lies below the estimate of $s$, by $2.5\cdot10^{-3}$ at $w = 3$ and
$6.7\cdot10^{-15}$ at $w = 25$ (Table~\ref{tab:odd} lists $w \ge 13$); on even cylinders the sector
$p = 0$ lies above it at every width we computed, but we have no argument that makes it a bound. The
comparison is with the estimate, so this is evidence, not a proof. We state the two sign properties
that the data show instead.

\begin{conjecture}\label{conj:Linf}
Let $a_0 = a_1 = 0$ and $a_k = \log\lambda_k$ for $k \ge 2$.
\begin{enumerate}
\item[(L$_\infty$)] $q \mapsto a_{2q}$ is convex: $a_{2q+2} - 2a_{2q} + a_{2q-2} \ge 0$ for all
  $q \ge 1$.
\item[(O)] $a_{k+1} - a_k \ge s$ for every odd $k$.
\end{enumerate}
\end{conjecture}

\begin{proposition}\label{prop:cond}
(L$_\infty$) implies $s \ge 0.565511911035$, and (O) implies $s \le 0.565511911049$.
\end{proposition}
\begin{proof}
By \S\ref{sec:model} and the proof of the lower bound of Theorem~\ref{thm:bounds},
$s = \lim a_k/k$. Under
(L$_\infty$) the increments $a_{2q+2} - a_{2q}$ are nondecreasing; since $a_{2q}/(2q) \to s$, they
increase to $2s$, so $s \ge (a_{24} - a_{22})/2$. Under (O), $s \le a_{24} - a_{23}$. The certificates
of \S\ref{sec:float} give, with outward rounding,
\begin{gather*}
  \log\lambda_{24} \ge 13.08824999346029, \qquad \log\lambda_{23} \ge 12.52273808241305,\\
  \log\lambda_{22} \le 11.95722617139012, \qquad \log\lambda_{24} \le 13.08824999346131;
\end{gather*}
the lower bounds are $24$ and $23$ times the certified per-site lower bounds at $w = 24, 23$ as
stored in \file{data/bounds.json} (before the outward rounding to $12$ decimals of
Table~\ref{tab:bounds}), minus $10^{-14}$ for their storage as floats, and the upper bounds come from Lemma~\ref{lem:cwup} with
$z = 0$. Hence $s \ge (13.08824999346029 - 11.95722617139012)/2 = 0.565511911035\ldots$ and
$s \le 13.08824999346131 - 12.52273808241305 = 0.565511911048\ldots$, rounded outward.
\end{proof}

Under both hypotheses $s$ is known to ten digits, $0.5655119110$. The estimate of
\S\ref{sec:value} lies $6.3\cdot10^{-12}$ above the lower and $7.7\cdot10^{-12}$ below the upper
conditional bound. Only the tails are used: the lower bound needs $a_{24} - a_{22} \le 2s$, which is
the Calkin--Wilf inequality that reflection positivity would give, and the upper bound needs (O)
at $k = 23$. The unconditional bounds remain those of Theorem~\ref{thm:bounds}.

\emph{Evidence.} Both statements hold wherever we can compute: the second differences of
$a_{2q}$ are positive for $2q = 2, 4, \dots, 24$, those of the odd widths are negative for
$k = 3, 5, \dots, 23$, and the increments $a_{k+1} - a_k - s$ alternate in sign for
$k = 0, \dots, 25$; at $k = 23$ the margin of (O) is $5.9\cdot10^{-12}$. One mode explains all of
this. The ratio of consecutive increments tends to $-0.338 = -e^{-1.083}$, the in-sector gap of
\S\ref{sec:z2}, so
\[
  g(k) = a_k - sk \approx -2\sigma + C(-\rho)^k, \qquad \rho = e^{-1.083},
\]
and a single sign, $C > 0$, gives both (L$_\infty$) and (O) for large $k$. A second, faster mode is
the splitting of the two parity sectors of the free strip, which alternates with period $4$ and
decays by a factor $-13.84$ per two columns (\S\ref{sec:model}). In the second differences of
$a_{2q}$ the two modes are $0.1147^q = \rho^{2q}$ and $(-0.072)^q$, and at $2q = 24$ the second is
$0.007$ of the first, shrinking by a factor $0.63$ per step.

The convexity is a quantitative property of unit weights, not a structural one. Its finite form
$Z(2q, m)^2 \le Z(2q-2, m)Z(2q+2, m)$ holds in all $1453$ cases we tested on the king's graph, but it
fails on the triangular lattice, for every unequal pair of diagonal weights we tried, and for equal
diagonal weights $x = \tfrac32, 2, 3$ on wide strips; along equal weights the two rates above cross,
on our data, between $x = \tfrac74$ and $x = 2$. The companion note \cite{rylko2026lnote} collects these tests
and the counterexamples, explains the onsets for nearly equal diagonal weights at $m = 4$ and $8$ by
the Krein signature of the left--right reflection, and lists the obstructions to the standard proof
techniques.

\begin{problem}
Prove (L$_\infty$) and (O) at unit weights, i.e.\ the sign $C > 0$ with a bound on the remainder
for every $k \ge 22$; Proposition~\ref{prop:cond} would then give ten rigorous digits. Another
route to ten digits is the sector bracket $f^{(1)}_w \le s \le f^{(0)}_w$ on even cylinders, which
holds in all our data (\S\ref{sec:model}) and has width $5.06\cdot10^{-11}$ at $w = 24$.
\end{problem}

\section{Reproducibility}\label{sec:repro}
All code and outputs are in the repository
\begin{center}
\url{https://github.com/dawidrylko/lattice-dimers},
\end{center}
tag \texttt{king-entropy-v1.0.0}, directory \texttt{research/king-entropy/}. The scripts need
Python~3.13 with numpy~2.5.3 \cite{numpy2020}, mpmath~1.4.1 \cite{mpmath} and scipy~1.18.1
\cite{scipy2020} (whose sparse eigensolver is ARPACK \cite{arpack1998}), as pinned in
\file{src/requirements.txt}. Every number in this paper is listed in \file{paper/NUMBERS.md},
with the committed output file it is read from and the command that wrote that file; the captions
and the text name the main ones. Many outputs record the commit of the code that wrote them. The
first version of this work is the tag \texttt{king-entropy-v0.1.0}, where the files sit directly in
\texttt{research/king-entropy/}; the commits recorded in the outputs of that version refer to its
history, which the tag keeps. The main commands, each run as \texttt{python3} followed by the script
and its arguments from \texttt{research/king-entropy/}:
\begin{center}\small
\begin{tabular}{ll}
\toprule
script & output \\
\midrule
\texttt{src/tm2.py king cyl 3 13}, \texttt{\dots\ 14 19} & \file{data/sec/sec\_king\_cyl\_*.json} \\
\texttt{src/cylw.py king cyl \textit{w} \textit{p}}, $w = 20, \dots, 25$ & \file{data/eig/eig\_king\_cyl\_w*\_p*.json} \\
\texttt{src/estimate.py} & \file{out/out\_estimate.txt}, \file{data/z2.json} \\
\texttt{src/lean.py king cyl \textit{w} \textit{p} longdouble} & \file{data/lean/lean\_*.json} \\
\texttt{src/errbar.py} & \file{out/out\_errbar.txt}, \file{data/errbar.json} \\
\texttt{src/bounds.py 24} & \file{out/out\_bounds.txt}, \file{data/bounds.json} \\
\texttt{src/runweight.py 3state \textit{w}\dots} & \file{out/out\_runweight\_3state.txt} \\
\texttt{src/audit\_bounds.py} & \file{out/out\_audit\_bounds.txt} \\
\texttt{src/pfaff\_rigor.py} & \file{out/out\_pfaff.txt} \\
\texttt{src/corner.py} & \file{out/out\_corner.txt}, \file{data/corner.json} \\
\texttt{src/surface.py} & \file{out/out\_surface.txt} \\
\texttt{src/cw.py} & \file{out/out\_cw.txt} \\
\texttt{src/zexact.py \dots}; \texttt{src/ltest.py} & \file{data/zex/zex\_*.json}, \file{out/out\_ltest.txt} \\
\texttt{src/lcert.py 25 30}, \texttt{\dots\ 26 30} & \file{out/out\_lcert\_m25.txt}, \dots \\
\texttt{src/lweight.py} [\texttt{big}] & \file{out/out\_lweight*.txt} \\
\texttt{L/symmetry.py}, \texttt{L/linf.py 16} & \file{L/out\_*.txt} \\
\texttt{L/involution.py}, \texttt{L/modes.py} & \file{L/out\_*.txt} \\
\texttt{src/validate.py}; \texttt{src/crosscheck.py} & \file{out/out\_validate.txt}, \dots \\
\texttt{paper/derived.py} & \file{paper/out\_derived.txt} \\
\bottomrule
\end{tabular}
\end{center}
The outputs were produced on two x86-64 machines with four cores and about $16$\,GiB of memory
each: the first version on one, the runs added in this version (extended precision, run weights,
the corner counts, the Pfaffian certificate and the two-mode fits) on the other. The extended-precision
runs use the 64-bit mantissa of the x87 \texttt{long double}, which \file{src/lean.py} asserts. The
algorithms are those of \S\ref{sec:model}; the implementation applies each cell step to reshaped
views of the state vector, without forming the transfer matrix. The widest
runs dominate the cost: the cylinder $w = 25$ ($8.4$\,GiB, $22.5$ min), the certified bounds at
$w = 24$ ($17$ min), the run-weighted bound at $w = 24$ ($13.3$\,GiB including the memory-mapped vector, $63$ min), the exact counts of the
$24 \times 24$ board ($1.5$\,GiB, $52$ min), and the certified (L) tests at $m = 26$ ($8.0$\,GiB, $43$ min). The script
\file{src/repro.py} re-runs \texttt{src/estimate.py}, \texttt{src/cw.py}, \texttt{src/bounds.py 16} and
\texttt{src/cylw.py king cyl 20 0} from a clean export of the repository and
checks that their outputs are identical up to timings and commit tags.

\section*{Data availability statement}
% TODO(submission): tag the release, archive it on Zenodo, then replace the placeholder.
The code and data that support the findings of this paper are openly available in the
repository \url{https://github.com/dawidrylko/lattice-dimers} at the tag
\texttt{king-entropy-v1.0.0}, and are archived on Zenodo,
DOI: TO-BE-ASSIGNED (Zenodo) \TODO{replace with the DOI minted for the tag}.

\section*{Disclosure statement}
The author reports no competing interests.

\bibliographystyle{amsplain}
\bibliography{refs}

\end{document}
